\documentclass[twocolumn]{aastex631}
\begin{document}
\title{A residual reionization ionizing-photon-budget shortfall for star-forming galaxies at z\~{}8 (lowxi systematic corner)}
\author{NebulaMind Lab (autonomous pipeline)}
\affiliation{NebulaMind Lab, \url{https://lab.nebulamind.net}}
\begin{abstract}
Reionization ionizing-photon-budget reconciliation at z\~{}8: star-forming galaxies require f\_esc=0.355 (+0.357/-0.181) to close the budget (Madau-Dickinson SFRD, log xi\_ion=25.5+/-0.15, clumping C=2-5, JWST-SFRD tail), versus indirect-proxy-inferred f\_esc=0.062 (+0.108/-0.039) from LzLCS O32/beta calibrations. Median delta(required-inferred)=+0.268 dex-frac (16-84\%: +0.069 to +0.626); 92\% of the systematic MC shows a shortfall. Conclusion: a genuine SHORTFALL remains, and the sign holds under both O32 and beta calibrations. The result is bounded by the xi\_ion x clumping x proxy-calibration systematic, not by statistics. Generated autonomously from public data (jwst) via the NebulaMind Lab runner: a bounded, reproducible, descriptive study.
\end{abstract}
\section{Introduction}
Recent studies on reionization have highlighted a potential crisis in the photon budget, suggesting that known sources may not be sufficient to account for the observed ionization [Muñoz2024]. This issue is further complicated by the increased demands on ionizing sources due to absorption-dominated reionization [Davies2021]. To address this problem, we need to reassess the contribution of star-forming galaxies to the ionizing photon budget. Previous efforts have focused on calibrating excursion set reionization models to conserve ionizing photons [Park2022] and evaluating the galaxy ionizing photon budget at high redshifts [Duncan2015]. However, a comprehensive understanding requires reconciling these findings with the analytic approaches to cosmic reionization [Madau2017].
\section{Data and method}
In this study, we adopt a literature-anchored approach, using the Madau \& Dickinson (2014) analytic fitting function for the cosmic star formation rate density (SFRD). We rely on published values for xi\_ion and O32/beta f\_esc proxy calibrations from LzLCS [Chisholm+22, Flury+22; Simmonds+24]. By focusing on systematics reconciliation over published literature values, we aim to provide a clearer picture of the ionizing photon budget during reionization.
\section{Result}
Our calculation reveals that star-forming galaxies require an escape fraction f\_esc = 0.355 (+0.357/-0.181) to close the reionization ionizing-photon-budget at z\~{}8. This value is significantly higher than the indirect-proxy-inferred f\_esc = 0.062 (+0.108/-0.039) derived from LzLCS O32/beta calibrations. The median difference between the required and inferred escape fractions is +0.268 dex-frac, with a range of +0.069 to +0.626. Notably, 92\% of our systematic Monte Carlo simulations show a shortfall in the photon budget.
\begin{figure}\centering\includegraphics[width=\columnwidth]{result.png}\caption{A residual reionization ionizing-photon-budget shortfall for star-forming galaxies at z\~{}8 (lowxi systematic corner)}\end{figure}
\section{Caveats}
It is essential to acknowledge that this study has limitations due to its reliance on automated, single-selection, and uncalibrated measurements. The accuracy of our results depends heavily on the assumptions made in the literature-anchored approach, particularly the Madau \& Dickinson (2014) SFRD fitting function and the adopted xi\_ion values. Additionally, the use of O32/beta f\_esc proxy calibrations introduces uncertainties that may not be fully captured by our systematic Monte Carlo simulations. Therefore, while our findings highlight a potential shortfall in the ionizing photon budget, further research is needed to refine these estimates and address the underlying systematics.
\section*{References}
\footnotesize
[Muoz2024] Muñoz et al. (2024). Reionization after JWST: a photon budget crisis?. bibcode 2024MNRAS.535L..37M \par [Davies2021] Davies et al. (2021). The Predicament of Absorption-dominated Reionization: Increased Demands on Ionizing Sources. bibcode 2021ApJ...918L..35D \par [Park2022] Park et al. (2022). Calibrating excursion set reionization models to approximately conserve ionizing photons. bibcode 2022MNRAS.517..192P \par [Duncan2015] Duncan et al. (2015). Powering reionization: assessing the galaxy ionizing photon budget at z < 10. bibcode 2015MNRAS.451.2030D \par [Madau2017] Madau (2017). Cosmic Reionization after Planck and before JWST: An Analytic Approach. bibcode 2017ApJ...851...50M

\end{document}
